\bibitem{auroux-local}
D. Auroux,
\emph{Mirror symmetry and T-duality in the complement of an anticanonical divisor},
J. G\"okova Geom. Topol. GGT \textbf{1} (2007), 51--91.
\url{https://arxiv.org/abs/0706.3207}.
Section 5.1, Proposition 5.2 の nodal model を参照。

\bibitem{cds-local}
A. Cannas da Silva,
\emph{Lectures on Symplectic Geometry}, Lecture Notes in Mathematics 1764, Springer, 2001;
revised January 2006.
\url{https://people.math.ethz.ch/~acannas/Papers/lsg.pdf}.
Sections 7.3, 9.2 の relative Moser と symplectic neighborhood theorem。

\bibitem{vungoc-local}
S. V\~u Ng\d{o}c,
\emph{On semi-global invariants for focus-focus singularities},
Topology \textbf{42} (2003), 365--380.
\url{https://arxiv.org/abs/math/0208245}.
Theorem 2.1、Proposition 3.1、Remark 3.2。

\bibitem{vw-local}
S. V\~u Ng\d{o}c and C. Wacheux,
\emph{Smooth normal forms for integrable Hamiltonian systems near a focus-focus singularity},
Acta Math. Vietnam. \textbf{38} (2013), 107--122.
\url{https://arxiv.org/abs/1103.3282}.
Theorems 1.2--1.4 は点近傍の smooth normal form とその automorphisms を扱う。

\bibitem{miranda-local}
E. Miranda and N. T. Zung,
\emph{Equivariant normal form for nondegenerate singular orbits of integrable Hamiltonian systems},
Ann. Sci. \'Ecole Norm. Sup. (4) \textbf{37} (2004), 819--839.
\url{https://arxiv.org/abs/math/0302287}.
Theorem 2.1 の orbit 近傍の normal form と、Section 2.3 の linear focus automorphism factor を参照。

\bibitem{auroux-blowup}
D. Auroux,
\emph{Special Lagrangian fibrations, wall-crossing, and mirror symmetry},
arXiv:0902.1595.
\url{https://arxiv.org/abs/0902.1595}.
Sections 3.1.2, 3.3.1 の blow-up model と torus-invariant model。

\bibitem{gross-toric}
M. Gross,
\emph{Examples of Special Lagrangian Fibrations},
arXiv:math/0012002.
\url{https://arxiv.org/abs/math/0012002}.
Sections 1--2 の Hamiltonian reduction と toric constructions。

\bibitem{cbm-conifold}
R. Casta\~no-Bernard and D. Matessi,
\emph{Conifold transitions via affine geometry and mirror symmetry},
Geometry \& Topology \textbf{18} (2014), 1769--1863.
\url{https://msp.org/gt/2014/18-3/gt-v18-n3-p12-p.pdf}.
Sections 3.4, 5.1 の explicit positive fibrations。

\bibitem{evans-gompf}
J. D. Evans,
\emph{Symplectic mapping class groups of some Stein and rational surfaces},
J. Symplectic Geom. \textbf{9} (2011), 45--82.
\url{https://arxiv.org/abs/0909.5622}.
Version 4, Section 5.2.1, Lemma 27 の局所 Gompf isotopy。

\bibitem{cg-at-blowup}
P. Chakravarthy and Y. Groman,
\emph{Almost toric fibrations on symplectic blow ups},
arXiv:2510.00994, version 1 (2025).
\url{https://arxiv.org/abs/2510.00994}.
Theorem 1.1, Definitions 2.2--2.3 の相対領域と divisor compatibility。

\bibitem{loi-zuddas-local}
A. Loi and F. Zuddas,
\emph{Symplectic maps of complex domains into complex space forms},
arXiv:0803.3532v1.
\url{https://arxiv.org/abs/0803.3532}.
Theorem 1.1 and the explicit radial maps in Section 2.


\bibitem{lerman-cut-local}
E. Lerman,
\emph{Symplectic cuts},
Math. Res. Lett. \textbf{2} (1995), 247--258.
\url{https://doi.org/10.4310/MRL.1995.v2.n3.a2}.
Section 1.1, Remarks 1.2 and 1.6 の cut と可換な Hamiltonian actions。

\bibitem{zung-focus-cf}
N. T. Zung,
\emph{A note on focus-focus singularities},
Differential Geometry and its Applications \textbf{7} (1997), 123--130.
\url{https://arxiv.org/abs/math/0110147}.
Section 3, Lemma 1 and Corollary 1 の focus monodromy。

\bibitem{efstathiou-martynchuk-cf}
K. Efstathiou and N. Martynchuk,
\emph{Monodromy of Hamiltonian systems with complexity 1 torus actions},
Journal of Geometry and Physics \textbf{115} (2017), 104--115.
\url{https://doi.org/10.1016/j.geomphys.2016.05.014}.
Theorems 3.2 and 3.15 の curvature と monodromy の関係。

\bibitem{cho-shin-maslov}
C.-H. Cho and H.-S. Shin,
\emph{Chern--Weil Maslov index and its orbifold analogue},
arXiv:1202.0556 (2012).
\url{https://arxiv.org/abs/1202.0556}.
Sections 2.3--2.4, Example 2.10 and Theorem 3.1 の符号と $2c_1$ の規約。


\bibitem{mcduff-uniruled-gc}
D. McDuff,
\emph{Hamiltonian $S^1$-manifolds are uniruled},
Duke Math. J. \textbf{146} (2009), 449--507.
\url{https://arxiv.org/abs/0706.0675}.
Corollary 1.2. 本文では負の単調性に必要な場合を orbit sphere から直接証明する。

\bibitem{gadbled-weights-gc}
A. Gadbled,
\emph{Families of monotone symplectic manifolds constructed via symplectic cut and their Lagrangian submanifolds},
Algebr. Geom. Topol. \textbf{11} (2011), 2319--2368.
\url{https://arxiv.org/abs/1001.2325}.
Appendix A.2 の equivariant clutching と weight sums。

\bibitem{achinger-lefschetz-gc}
P. Achinger,
\emph{Fundamental groups in algebraic geometry},
lecture notes, 19 March 2025.
\url{https://achinger.impan.pl/pi1/notes.pdf}.
Section 2.3, Corollary 2.3.4 の smooth complete intersections の単連結性。


\bibitem{circle-morse-ncb}
K. A. de Rezende, G. G. E. Ledesma, O. Manzoli Neto and G. M. Vago,
\emph{Lyapunov graphs for circle valued functions},
Topology Appl. \textbf{245} (2018), 62--91.
\url{https://doi.org/10.1016/j.topol.2018.06.008}.
Theorem 5.5, Lemma 5.7 and Theorem 5.9 の circle-valued Morse functions と graph。

\bibitem{hatcher-bct}
A. Hatcher,
\emph{Algebraic Topology}, Cambridge University Press (2002).
\url{https://pi.math.cornell.edu/~hatcher/AT/AT.pdf}.
Section 1.3 の universal covering と lifting criterion，
Propositions 0.18--0.19 の homotopy pushout，
Theorems 4.32, 4.37, 4.41, 4.57 と Example 4.45 の
Hurewicz theorem，relative Hurewiczの基本群作用によるquotient，
homotopy長完全列，cohomologyの $K(G,n)$ 表現，Hopf bundle。


\bibitem{bel-fractional-tcf}
H. Broer, K. Efstathiou and O. Lukina,
\emph{A geometric fractional monodromy theorem},
Discrete Contin. Dyn. Syst. Ser. S \textbf{3} (2010), 517--532.
\url{https://doi.org/10.3934/dcdss.2010.3.517}.
Section 2.1, Remark 2 の $\mathbb Z/2$ isotropy と reverse hyperbolic 型。
本稿の weighted model に対する global fractional monodromy は主張しない。


\bibitem{randal-williams-hypersurface}
O. Randal-Williams,
\emph{Monodromy and mapping class groups of 3-dimensional hypersurfaces},
Math. Ann. \textbf{391} (2025), 1965--2003.
\url{https://doi.org/10.1007/s00208-024-02951-4}.
Section 2.1: smooth degree-$d$ threefold hypersurface の整数cohomology、
$H^3$ の自由性およびChern類からのEuler数．


\bibitem{akbulut-double-fishtail}
S. Akbulut,
\emph{Double knot surgeries to $S^4$ and $S^2\times S^2$},
arXiv:1103.3137v2 (2011).
\url{https://arxiv.org/abs/1103.3137}.
Introduction and Figure 1: $S^4$ の二つの fishtails への既知の分解を記述する．
本稿の actual 埋込みとの微分同相や symplectic 同一視は仮定しない．

\bibitem{asimov-round}
D. Asimov,
\emph{Round handles and non-singular Morse--Smale flows},
Ann. of Math. (2) \textbf{102} (1975), no.~1, 41--54.
\url{https://doi.org/10.2307/1970972}.
Round handleの標準概念を参照．本稿で用いる分解はHopf商から直接証明する．

\bibitem{cieliebak-eliashberg-flexible-local}
K. Cieliebak and Y. Eliashberg,
\emph{Flexible Weinstein manifolds}, arXiv:1305.1635 (2013).
\url{https://arxiv.org/abs/1305.1635}.
Section 1.1, pp.~2--3 の Weinstein cobordism とisotropic stable manifoldsによるindex制限．

\bibitem{brown-relative-pi2}
R. Brown,
\emph{On the second relative homotopy group of an adjunction space:
an exposition of a theorem of J.~H.~C. Whitehead},
J. London Math. Soc. (2) \textbf{22} (1980), 146--152.
\url{https://doi.org/10.1112/jlms/s2-22.1.146}.
Author PDF: \url{https://www.groupoids.org.uk/pdffiles/whiteheadthm.pdf}.
Section 2 の relative crossed module と Peiffer identity．

\bibitem{conrad-group-extensions}
B. Conrad,
\emph{Math 210B. Group cohomology and group extensions},
course handout, Stanford University.
\url{https://math.stanford.edu/~conrad/210BPage/handouts/gpext.pdf}.
Section 2 の $H^2$ によるcentral extensionの分類．
